Figure~\ref{fig:affected-mse} summarizes the original row-index development
evaluation. These preserved figures are historical context, not the primary
grouped estimate.

\begin{figure}[t]
    \centering
    \includegraphics[width=\linewidth]{v02_affected_mse_bar.png}
    \caption{Affected-region MSE for identity, threshold, Thayer-Direct,
    Thayer-Residual, Thayer-BR v0.1, and Thayer-BR v0.2 Moderate.
    Affected-region metrics focus evaluation on the thresholded blend-to-target
    change region. Values are original development results.}
    \label{fig:affected-mse}
\end{figure}

\begin{figure}[t]
    \centering
    \includegraphics[width=\linewidth]{v02_improvement_ratio.png}
    \caption{Original development improvement ratio versus identity. These
    historical values are retained separately from grouped development results.}
    \label{fig:normal-stress-improvement}
\end{figure}

On the original 1,000-blend comparable normal development evaluation, identity
affected-region MSE is 0.068122. Thayer-Direct reaches 0.004236,
Thayer-Residual reaches 0.004431, Thayer-BR v0.1 reaches 0.002451, and
Thayer-BR v0.2 Moderate reaches 0.002108. This corresponds to about 32.3x
lower affected-region MSE than identity for v0.2 Moderate.

On the hard stress test, identity affected-region MSE is 0.075541.
Thayer-Direct reaches 0.009390, Thayer-Residual reaches 0.007069,
Thayer-BR v0.1 reaches 0.004587, and Thayer-BR v0.2 Moderate reaches 0.003847.
This corresponds to about 19.6x lower affected-region MSE than identity for
v0.2 Moderate.

The source-group-disjoint v0.2 retrain gives 0.00231890 affected MSE on grouped
normal development data (28.8127x lower than identity), 0.00458983 on hard
stress (15.8025x), 0.00872771 on compact-bright stress (9.18304x), and
0.00491680 on high-core-obstruction stress (15.8378x). Worse-than-identity
counts are 0, 3, 2, and 1 out of 1,000, respectively. These grouped values are
the more defensible development estimates, but remain non-final and come from
one training seed.

\begin{figure}[t]
    \centering
    \includegraphics[width=\linewidth]{v02_core_mse.png}
    \caption{Original row-index development core affected MSE comparison.
    Thayer-BR v0.2 Moderate improves stress core affected MSE from 0.013848 for
    Thayer-BR v0.1 to 0.009533.}
    \label{fig:core-mse}
\end{figure}

\begin{figure}[t]
    \centering
    \includegraphics[width=\linewidth]{v02_weighted_vs_v01_scatter.png}
    \caption{Original row-index development per-sample affected-region MSE for
    Thayer-BR v0.2 Moderate versus Thayer-BR v0.1. Points below the parity line
    favor v0.2 Moderate.}
    \label{fig:v02-v01-scatter}
\end{figure}

On the original row-index development suites, compared with Thayer-BR v0.1,
Thayer-BR v0.2 Moderate lowers normal
affected-region MSE by about 14\%, stress affected-region MSE by about 16\%,
and stress core affected MSE by about 31\%. Normal worse-than-identity cases
fall from 1/1000 for v0.1 to 0/1000 for v0.2 Moderate; both models have 0/1000
stress worse-than-identity cases.

\begin{figure}[t]
    \centering
    \includegraphics[width=0.75\linewidth]{v02_multiseed_summary.png}
    \caption{Evaluation-seed affected-MSE ratio for the original Thayer-BR v0.2
    Moderate checkpoint. Across
    three normal and three stress evaluation seeds, the model reaches
    32.02 +/- 1.21x normal and 19.55 +/- 0.30x stress. These are not independent
    training seeds.}
    \label{fig:multiseed}
\end{figure}

In the original development ablation, the stronger v0.2 weighting slightly
improves stress core MSE
relative to Moderate, but worsens aggregate normal affected MSE, stress
affected MSE, and stress non-core affected MSE. Moderate weighting is therefore
the best current v0.2 setting.

\begin{figure}[t]
    \centering
    \includegraphics[width=\linewidth]{v02_weighted_improvement_example.png}
    \caption{Qualitative example where Thayer-BR v0.2 Moderate improves over
    Thayer-BR v0.1.}
    \label{fig:v02-success}
\end{figure}

\begin{figure}[t]
    \centering
    \includegraphics[width=\linewidth]{v02_counterexample.png}
    \caption{Qualitative counterexample where Thayer-BR v0.1 beats
    Thayer-BR v0.2 Moderate on an individual sample. Such cases motivate
    visual-vs-metric diagnostics and limitation figures.}
    \label{fig:v02-counterexample}
\end{figure}
